\scanpage{49}
% loc: cancelled continuation at head
\deleted{regarder $V$ seulement (où $\overline V=V\otimes_{\mathbb Q}\overline{\mathbb Q}$)}
\[
\deleted{\overline V^{G}\subset\overline V\cap g(\overline V)}
\]
\deleted{et de prouver que c'est une identité [car $\overline V^G\cap V=V^G$].}

% loc: first equivalence
Or la relation $x\in\overset{\deleted{\unsure{-}}}{V}\cap g(\overline V)$ équivaut :
\begingroup\makeatletter\tagsleft@true\makeatother
\[
x\in\overset{\deleted{-}}{V},\quad g^{-1}x\in\overline V \tag*{\unsure{ds|$\alpha$)}}
\]
\endgroup
% loc: locus and second equivalence
et implique évidemment, si $\unsure{\overline T|\overline J}$\inserted{$\subset\overline G$} est \unsure{le|la} \deleted{\unsure{déclinaison|clôture} de} \inserted{lieu de} $g$ sur $G_{\overline{\mathbb Q}}$
\[
x\in\overline V,\quad\deleted{\unsure{g^{-1}x}}\quad t\deleted{\unsure{^{-1}}}s^{-1}\in G_x(\mathbb C)\quad\text{pour }t,s\in\overline T(\mathbb C),
\]
et $=$ y équivaut.
% loc: stabilizer paragraph
En d'autres termes, si $G_x$ est le stabilisateur de \deleted{$\unsure{G}$}$x$, qui est un ss-groupe algébrique de $G$, \deleted{\unsure{on}} \inserted{la relation $x\in V\cap g(\unsure{\overline W|\overline V})$ équi-}vaut :
\[
\deleted{\unsure{g}}\overline R\unsure{\in|\subset}\overline{G_x}\deleted{\unsure{(\mathbb C)}},
\]
où $\overline R$ est l'adhérence de l'im. $\overline T\times\overline T\longrightarrow\overline G$, $(s,t)\longmapsto s^{-1}t$.

% loc: image subgroup and equality question
Or soit $\widehat R$ l'image de $\overline R$ dans $G$, $H$ le sous-gr. qu'elle engendre, $H\subset G$. \deleted{Donc \unsure{la} propriété souhaitée équivaut à quelle \unsure{?}} La relation ds équivaut donc à $\deleted{\unsure{R}}H\subset G_x$ i.e. $H$ invarie $x$.
% loc: left insertion at opening of proposition
\inserted{La propriété \unsure{qu'on|qui} \unsure{veut est|nous est} celle-ci :}
% loc: heavily cancelled proposition and replacement
Pour tt $G$-module $V$, \deleted{et tt $x\in V$, tel que le stabilisateur $G_x$ contienne $\unsure{R}$, \inserted{i.e. tel que $R$ invarie $x$,} on a que $x\in V^G$ i.e. $G_x=G$.}
\deleted{Il suffit pour \unsure{ceci} que \unsure{le} \illegible{1} soit le ss-gpe engendré par $R$, soit égal à $G$. C'est aussi \unsure{nécessaire}} que les invariants de $H$ et de $G$ (\deleted{pour} \unsure{tt} repr. de $G$) soient les $=$\unsure{s}. Il suffit que l'on ait $H=G$, mais cette condition n'est pas nécessaire à priori, car un $k$-groupe de Borel de $G$ a cette propriété.
% loc: NB at left of cancelled lines
\margin{NB $\deleted{\unsure{\overline T}}\overline R,\overline H$ ne dépendent pas du choix de $g$ mod $G(\overline{\mathbb Q})$, \deleted{\unsure{à un près|a priori}}}
% loc: sloping note above 5
% slanted
\margin{\uline{Regarder} motifs \uline{associés} à \uline{courbes elliptiques}}
% loc: item 5, first half
5) Bien entendu, sans $=$ supposer que $k$ est alg. sur $\mathbb Q$, on peut regarder le foncteur (4.1). Mais il n'est plus en général pl. fidèle, $=$ si $k$ est de t.f. sur $\mathbb Q$, car pour $\unsure{\mathcal M|M}$ \inserted{$\otimes$-cat.} fixée de type fini $\subset\operatorname{Mot}(k)$, à condition d'agrandir $k$, on trouve que \inserted{les} $T_{\mathrm{DR}}(M)$ \uline{contiennent} $T_{\mathrm B}(M)$, donc les hom. des structures mixtes sont simplt les hom entre les $T_{\mathrm B}(M)$ : idiot !
% loc: closing lines of item 5
Pour $\mathcal M$ fixé, on peut espérer tout au mieux avoir quelque résultat de \unsure{disjonction} de $T_{\mathrm B}$ et $T_{\mathrm{DR}}$ si $k$ est $\pm$ \unsure{modeste|modérée|modulaire} (p.ex avoir des \unsure{estimations} de degré de transcendance minimum).
% loc: paper watermark, read sideways and compared across sheets
% printed form
% paper watermark, visible letters in the direction of the paper lettering
\par\noindent\textsc{Extra Strong}\\
\unsure{Maunoury|Maunouy} 1808
% archive sheet number: 48

